\section{Priced logarithmic steps}\label{sec:engine}

This section constructs the optimization primitive used by the matching
algorithm and the local projection procedure.  The primitive maintains
exact values implicitly, while exposing sufficiently accurate coordinate
registers to its client.  Its progress guarantee depends on the length of
the actual path representation returned by a cycle data structure.  In
particular, cancellation in the resulting circulation does not reduce the
motion charged to a step or to a register refresh.

Approximate minimum-ratio circulations serve as optimization directions
in the almost-linear flow framework of Chen et al.~\cite{FlowRandomized}.
The step estimate below is proved for the logarithmic objective and the
raw-motion convention used here.

\subsection{The interface}\label{engine:interface}

Fix a master bound $S=2^s$ and integer parameters $a\ge1000$ and
$\sigma\ge0$.  Recall
\begin{equation}\label{engine:parameters}
 B=2^{\lceil s/a\rceil},\qquad q=B^{-100},\qquad
 \delta=S^{-1000},\qquad \varepsilon=S^{-3000}.
\end{equation}
Every asymptotic estimate in this section is for fixed parameters.  The
constants multiplying $1/a$ in exponents are absolute; an $o_a(1)$ term
may include logarithmic powers whose exponents depend on $a$ and
$\sigma$.  We use deterministic prerequisite checks, including
$q\ge S^{-1}$.  An inadequate prerequisite suspends the parameter branch
without producing an answer.  A detected stochastic failure is a distinct
outcome, with the probability bound stated below.

Let a changing oriented multigraph have real edge variables $f_e$ on the
grid $\varepsilon\mathbb Z$.  Loops and individual parallel edge copies
are allowed.  Its separable objective $\Phi(f)=\sum_e\phi_e(f_e)$ has one
of the following terms on each edge:
\begin{equation}\label{engine:objectives}
 \log f,\qquad
 \log(f-D_1)+\log(D_2-f),\qquad
 -W\log(T-f),\quad 1\le W\le S^{20},
\end{equation}
where $W$ is an integer.  The endpoints $D_1,D_2,T$ lie on the same grid
as the values.  Write $g_e=\phi'_e(f_e)$, and let $\lambda_e$ be the sum
of the absolute derivatives of the individual logarithmic terms on that
edge, including their weights.  Thus, for the two-sided term,
\[
 g_e=\frac1{f_e-D_1}-\frac1{D_2-f_e},\qquad
 \lambda_e=\frac1{f_e-D_1}+\frac1{D_2-f_e}.
\]
All data and values have magnitude at most $S^{20}$, and the positive
arguments at pricing and refresh states lie between $\delta/10$ and
$S^{20}$.  Small additional relative margins immediately after a step
are harmless in these bounds.  The applications maintain the stricter
pre-step margins established in the preceding sections.

The client may add, remove, or retarget edges and vertices, change a
descriptor or value on one coordinate, and request an exact point value.
Its explicit value settings need not be circulations.  Between such
edits, all increments performed by the primitive must be circulations.
The values themselves need not have zero divergence.

\begin{lemma}[Priced logarithmic steps]\label{lem:priced-steps}
Suppose a nominal bound $R$ counts the logical vertex and edge births and
endpoint insertions over one instance, including retargetings.  Mere
value or price changes on retained topology need not be counted as new
incidences.  Assume
\begin{equation}\label{engine:nominal-size}
                         R+B\le S^{1+2/a}.
\end{equation}
There is a uniform implementation with the following properties.
\begin{enumerate}[label=\textup{(\roman*)}]
\item A pricing call may return \textnormal{\textsc{Low}}, but it cannot do so if
there is a nonzero circulation $d$ satisfying
\begin{equation}\label{engine:high-direction}
                    g^{\mathsf T}d\ge\tfrac12\|\lambda d\|_1,
\end{equation}
unless it reports a detected failure or suspends.
\item A call not returning \textnormal{\textsc{Low}}, failure, or suspension performs
an exact increment $h\in\varepsilon\mathbb Z^E$ that is a circulation and
satisfies
\begin{equation}\label{engine:step-guarantee}
       \|\lambda h\|_1\le q/100,
       \qquad g^{\mathsf T}h\ge q^2/1000.
\end{equation}
It leaves all logarithmic arguments positive and increases $\Phi$ by
$\Omega(q^2)$.
\item The implementation exposes registers $\widehat f_e$.  After
settling updates and before the next pricing call, the descriptors are
valid at these registers and
\begin{equation}\label{engine:register-guarantee}
       \lambda_e(\widehat f_e)|f_e-\widehat f_e|\le q/100.
\end{equation}
An explicit client setting can reanchor the coordinate exactly.
Apart from such client settings, the number of refresh notifications is
at most the number of successful pricing steps times
$S^{O(1/a)+o_a(1)}$.  Notifications are supplied as a list of dirty
coordinates.  Maintaining the representation without taking a step
causes no implicit primal motion.
\item With $Q$ external pricing, point, and edge operations, total time,
excluding the client's own work, is
\begin{equation}\label{engine:time}
                     (R+B+Q)S^{O(1/a)+o_a(1)}.
\end{equation}
The bound is on every computation path and includes initialization,
arithmetic, random bits, and all representation changes.  Polynomially
many operations are allowed, with a loose guard $Q\le S^{100}$.
\end{enumerate}
The conditional probability of detected failure in the entire instance
is at most $2^{-\sigma}S^{-60}$, including when the instance is chosen
adaptively.  Deterministic size, precision, and representation checks may
suspend, but none does so for all sufficiently large $S$ at each
sufficiently large fixed $a$ and each fixed $\sigma\ge0$ within the
stated hypotheses.
\end{lemma}

The objective conclusion in part~\textup{(ii)} follows directly from its
two inequalities once positivity is established.  The sum of absolute
relative argument changes is at most $\|\lambda h\|_1$.  For a positive
logarithm, $\log(1+z)\ge z-z^2$ when $|z|\le q/100$; the sum of these
squared relative changes is at most $q^2/10000$.  A negative logarithm
in~\eqref{engine:objectives} is convex and lies above its supporting
tangent.  Hence the total increase is at least
$q^2/1000-q^2/10000$.

\subsection{A backend with explicit forest paths}\label{engine:backend-interface}

We use the following fully dynamic backend.  Its construction, rather
than a dynamic optimization oracle, is given in
Sections~\ref{sec:backend-quality} and~\ref{sec:backend-updates}.  The
minimum-ratio-cycle and hierarchical-routing viewpoint is related to
the flat-forest cycle machinery of~\cite{IncrementalGraphs}.  We give a
complete construction for the precise interface below, including its
raw path length and explicit forest recourse guarantees.

\begin{lemma}[Flat-cycle backend]\label{lem:flat-cycle}
Let $M=2^u$, where $u$ is divisible by $2a^2$, and suppose
\begin{equation}\label{engine:backend-size-range}
                         S^{1/a}\le M\le S^2.
\end{equation}
Consider a dynamic undirected graph with reference orientations and
fixed vertex slots, of total current size at most $M$ and maximum degree
at most three.  It has positive integral lengths $\ell_e$ and integral
signed gradients $\gamma_e$, both of magnitude at most $M^{1200a}$.
Insertions and deletions are arbitrary; a price replacement is a
deletion followed by an insertion.  Put
\[
 \OPT=\max_{d\ne0,\ d\text{ a circulation}}
          \frac{|\gamma^{\mathsf T}d|}{\|\ell d\|_1},
\]
with value zero if there is no nonzero circulation.

The backend exposes a forest.  Every forest vertex maps to an input
vertex slot; every forest edge either maps to a current input edge with
specified endpoint correspondence, or is null and joins equal images.
The multiplicity over each input slot is at most $M^{10/a}/4$.
A query returns a list of oriented single input edges and oriented
forest paths, the latter specified by endpoints.  Each piece is used
with unit multiplicity, with repeated pieces listed separately.  Their
signed physical projection is a circulation $c$.  Its \emph{raw length}
$V$ is the sum of input lengths over all mapped edge traversals with
repetition; a null forest traversal contributes zero.  If $\OPT>0$, the
answer has $V>0$ and
\begin{equation}\label{engine:flat-quality}
                   \frac{|\gamma^{\mathsf T}c|}{V}
                       \ge M^{-10/a}\OPT.
\end{equation}
An empty or missing answer is permitted when $\OPT=0$.
There are $M^{O(1/a)+o_a(1)}$ pieces in an answer.

Initialization costs $M^{1+O(1/a)+o_a(1)}$.  Amortized edit and query
costs, including explicitly exposed forest logs, are
$M^{O(1/a)+o_a(1)}$, with bounds on every path.  Retained forest
occurrences and edges have stable lifetimes.  Endpoint or mapping changes
are replacements, and removals are exposed before the old mappings or
input data are invalidated.  Transactions permit deletion-first
old-to-final updates with no intervening path query.

The bounds apply over the polynomial lifespans required by
Lemma~\ref{lem:priced-steps}, and against adaptive clients.  Reliability
can be chosen so that detected failure in an entire engine instance has
conditional probability at most $2^{-\sigma}S^{-60}$.  A flagged state
has no quality requirement.  Deterministic prerequisite checks may
suspend, with eventual adequacy as in Lemma~\ref{lem:priced-steps}.
\end{lemma}

\subsection{Degree reduction and integer pricing}\label{engine:pricing}

For each logical edge insertion, create a fresh port at each endpoint;
a loop receives two distinct ports.  The ports belonging to one logical
vertex lifetime form a path, extended on further endpoint insertions.
Retargeting uses fresh ports.  A retired port may remain on its path, and
a retired vertex may leave a wholly inactive path after all its real
incidences have disappeared.  Distinct vertex lifetimes have distinct
paths.  Mere value or price changes retain the ports.  Give every path
connector length one and gradient zero.  The resulting graph has maximum
degree three and linear lifetime size in $R$.

Choose $M$ by rounding $16(R+B)$ upward to a power of two whose exponent
is divisible by $2a^2$.  For fixed $a$ this introduces only a fixed
factor.  Reserve, for example, $M/2$ vertex slots; the ports and any
isolated logical-vertex representatives fit these slots, and slots plus
edges fit the backend bound.  These counts and
\eqref{engine:backend-size-range} are checked.  They hold eventually by
\eqref{engine:nominal-size}.  Identifiers can be renumbered using
dictionaries.

Summing conservation over one port path cancels its connectors and
recovers conservation at the logical vertex.  Conversely, a logical
circulation $d$ lifts to the port graph by routing its incidence
imbalances along each path.  The total absolute connector coefficient is
at most
\begin{equation}\label{engine:connector-cost}
                   2M\sum_{e\text{ real}}|d_e|.
\end{equation}
Indeed each real incidence contributes its coefficient to at most $M$
connector positions, and each real edge has two incidences.  No connector
primal value needs to be stored; connectors represent conservation of
increments only.

An anchor of coordinate $e$ records its exact current value
$\widehat f_e$.  Set $H=S^{100}$ and price real edges by
\begin{equation}\label{engine:integer-prices}
  \ell_e=\bigl\lceil H\lambda_e(\widehat f_e)\bigr\rceil,
  \qquad
  \gamma_e=\operatorname{round}\bigl(Hg_e(\widehat f_e)\bigr).
\end{equation}
The sum or difference of at most two reciprocals of positive dyadic
arguments, with the allowed integer weight, can be priced by exact
bounded-length integer quotient arithmetic.  No logarithm is evaluated.
At anchors,
\begin{equation}\label{engine:derivative-range}
             S^{-20}\le\lambda_e(\widehat f_e)\le S^{1022}.
\end{equation}
The small post-step margins do not affect this loose upper estimate.
Since $S\le M^a$, these prices fit the backend range $M^{1200a}$.

For the moment assume the register estimate
\eqref{engine:register-guarantee}.  Put $\tau=q/100$.  An individual log
argument $r$ with weight $W\ge1$ changes by a relative amount at most
$\tau/W\le\tau$ between an anchor and a pricing call.  Its derivative
magnitude changes by at most $\tau/(1-\tau)$ times its anchored
magnitude.  Summing these errors term by term remains valid when the
signed derivatives nearly cancel.  Together with rounding error at most
$1/(2H)$ and $\ell_e\ge S^{80}$, this gives, at every pricing call,
\begin{equation}\label{engine:price-accuracy}
 \left|g_e(f_e)-\frac{\gamma_e}{H}\right|
       \le\frac q{10}\frac{\ell_e}{H},\qquad
  .9\frac{\ell_e}{H}\le\lambda_e(f_e)
       \le1.1\frac{\ell_e}{H}.
\end{equation}
The guard $q\ge S^{-1}$ makes the integer rounding error negligible
uniformly in the coordinate.  These inequalities require no scan of
the exact current vector: anchors are checked locally, and the
displacement invariant below supplies the intervening estimate.

If~\eqref{engine:high-direction} holds, its lift has integer-priced ratio
at least $1/4$.  To see this, the numerator in units of $H$ is at least
\[
 \tfrac12\|\lambda d\|_1
       -\frac q{10}\sum_{e\text{ real}}\frac{\ell_e}{H}|d_e|,
\]
whereas~\eqref{engine:price-accuracy} compares the real length sum with
$\|\lambda d\|_1$.  The additional connector length in
\eqref{engine:connector-cost} is negligible, since
$2M\le2S^2$ while every real length is at least $S^{80}$.
The bound $1/4$ follows with room to spare for sufficiently large checked
sizes.  We also check
\begin{equation}\label{engine:quality-margin}
                         M^{-10/a}/4>2q.
\end{equation}
It holds eventually: $M\le S^2$, whereas
$q=2^{-100\lceil s/a\rceil}\le S^{-100/a}$.

Compute the exact gradient sum and raw length of a backend answer using
the path aggregates constructed below.  Accept it precisely when
\begin{equation}\label{engine:accept}
                         V>0,\qquad
                  |\gamma^{\mathsf T}c|\ge2qV;
\end{equation}
otherwise return \textsc{Low}.  Lemma~\ref{lem:flat-cycle} and
\eqref{engine:quality-margin} show that this cannot miss
\eqref{engine:high-direction} in an unfailed proceeding computation.

Reverse every piece of an accepted answer if necessary so that
$\gamma^{\mathsf T}c>0$, and delete its connector coordinates from the
logical increment.  If $N_e$ counts every real
traversal of $e$ in its pieces, then $|c_e|\le N_e$.  In particular,
\eqref{engine:price-accuracy} and~\eqref{engine:accept} imply
\begin{equation}\label{engine:accepted-direction}
       g^{\mathsf T}c\ge qV/H,\qquad
       \sum_{e\text{ real}}\lambda_e N_e\le1.1V/H.
\end{equation}
The first bound could be strengthened to $1.9qV/H$; the weaker version
is convenient.  With $c$ now oriented for ascent, choose a positive power
of two $\mu$ with
\begin{equation}\label{engine:step-scale}
                     q/400\le \mu V/H\le q/200.
\end{equation}
The interval has ratio two and therefore contains such a power.
Applying $\mu$ along all pieces gives $h=\mu c$, with
\[
          g^{\mathsf T}h\ge q^2/400,
          \qquad \|\lambda h\|_1\le1.1q/200.
\]
These are stronger than~\eqref{engine:step-guarantee}.
Moreover, the second inequality of
\eqref{engine:accepted-direction} bounds the sum of absolute partial
motions during the entire list of traversals, not just its net motion.
Every individual argument consequently stays within relative distance
$1.1q/200$ of its starting value, and within that distance divided by its
weight for a weighted term.  All arguments remain positive.  The whole
step is applied before repricing or servicing client changes; individual
pieces need not themselves be circulations.  Exact grid compatibility
of $\mu$ is verified below.

\subsection{Oriented forest operations}\label{engine:dynamic-trees}

We implement path sums, path increments, and threshold detection with
link-cut trees~\cite{DynamicTrees}.  Signed path updates, separate absolute
accumulation, and length-weighted threshold reporting are the dynamic-tree
operations in~\cite[Lemma~2.7]{IncrementalGraphs}.  Its flow-maintenance
construction also transfers implicit forest contributions into explicit
edge values when mapped edges change~\cite[Theorem~3.6]{IncrementalGraphs}.
For our interface, the signed-value and absolute-motion invariants must
survive dormant-copy reuse and repricing transactions.  We give the
augmentation and amortization here, then specify the copy transfers and
reset ordering in Sections~\ref{engine:exact-values} and~\ref{engine:refresh}.

Subdivide each forest edge by an edge node.  The other nodes are
occurrence vertices, and represented roots are always occurrence
vertices.  On an edge node $c$ mapping to an oriented actual edge, let
$s_c=+1$ when its represented-parent side is its tail and $s_c=-1$ when
that side is its head.  A null edge may receive either reference
orientation.  A real-variable copy stores two different accumulations:
\begin{itemize}
\item $h_c$ is signed accumulation since its creation, in the fixed
orientation of the mapped actual edge;
\item $A_c$ is absolute accumulation since the coordinate's most recent
anchor or the copy's creation, whichever occurred later.
\end{itemize}
After everting an occurrence vertex $v$ and exposing the path to $w$,
the signs $s_c$ describe traversal from $v$ to $w$.  A path increment
uses
\begin{equation}\label{engine:path-add}
                    h_c\leftarrow h_c+\mu s_c,
                    \qquad A_c\leftarrow A_c+|\mu|
\end{equation}
on real-variable copies.  Connector and null copies do not need these
primal accumulations.

Preferred-path aggregates store the signed sum
$\sum_c s_c\gamma_c$ and the raw sum $\sum_c\ell_c$ on mapped nonnull
edges.  They also store a witnessed minimum of $\theta_c-A_c$ over
enabled real copies, where the thresholds are specified in
Section~\ref{engine:refresh}.  Ineligible nodes are omitted from this
minimum.  Disabled real copies still receive the updates in
\eqref{engine:path-add}.

Evert flips signs on edge nodes of the reversed exposed path, but does
not negate their stored $h_c$.  The reversal, signed-addition, and
absolute-addition tags have the constant-size algebra
\begin{equation}\label{engine:tag-algebra}
 (s_c,h_c,A_c)\longmapsto
       (r s_c,\ h_c+\alpha s_c,\ A_c+\beta),\qquad r\in\{-1,+1\},
\end{equation}
together with reversal of path order when appropriate.  Applying one
such transformation and then another composes their signed offsets as
$\alpha_1+r_1\alpha_2$.  Thus a later reroot does not retroactively
negate an earlier physical increment.  The gradient aggregate changes
sign on reversal, the raw length does not, and the enabled minimum
shifts by the absolute-addition tag with its witness retained.  Tags are
pushed before modifying children or point fields.  Only preferred-path
aggregates are combined; no virtual-subtree aggregate is used for these
operations.

Access at a node does not change any represented-parent direction; it
only changes which represented edges are preferred.  Thus it does not
alter the signs $s_c$.  Point accesses without evert recover exact
individual data.  To insert an edge node, link it as a new child on its
tail side, initializing sign $+1$; evert the head endpoint in its
separate component and link that endpoint below the edge node.  To
delete an edge node, inspect its sign by point access, cut its child
endpoint from its parent, and then cut the edge node from its parent.
No evert at an edge node is needed.  Forest transactions give deletions
before invalidation and insertions into the final forest, so every link
joins different components.

\begin{lemma}[Augmented forest cost]\label{engine:tree-cost}
On a forest of at most $N_f$ subdivided nodes, the preceding operations
use $O(\log N_f)$ amortized tree primitives per link, cut, evert, expose,
point operation, or lazy path update.  Reporting and disabling a
witnessed threshold violation costs the same amortized amount per
reported node.  Each primitive uses a constant number of operations on
the stored fields.
\end{lemma}

\begin{proof}
Use auxiliary splay trees in down-path order for preferred
paths~\cite[Section~6]{SplayTrees1985}. Their
roots keep path-parent pointers to the represented parents of their path
tops.  Access at $v$ starts at $v$ and repeatedly splays the current node
$y$, replaces its right subtree by the path exposed so far, and advances
to its previous path parent.  The demoted path receives path parent $y$.
A final splay of $v$ exposes the represented-root-to-$v$ path.  Evert
reverses that path.  A cut from parent accesses its endpoint and detaches
the left ancestor subtree.  A link accesses its endpoints and installs a
path-parent pointer for a represented root.  The stated tags require
only bounded work during each pointer or aggregate change.

For the amortization, give a node weight one plus the sum of represented
subtree sizes of its nonpreferred children.  Let $s(x)$ be its auxiliary
subtree weight and $r(x)=\log_2 s(x)$, and use
$\Phi=\sum_x r(x)$.  Node weights stay fixed during a splay.  The usual
rotation inequalities from the access lemma~\cite[Lemma~1]{SplayTrees1985}
can be recalled briefly.  In a double rotation
with node $x$, parent $y$, and grandparent $z$, we have $r'(x)=r(z)$ and
$r(y)\ge r(x)$.  For zig-zag,
$r'(y)+r'(z)\le2r'(x)-2$ by disjoint child weights.  For zig-zig,
$r(x)+r'(z)\le2r'(x)-2$ by the disjoint old-$x$ and new-$z$ subtrees,
and $r'(y)\le r'(x)$.  Thus two rotations plus their change in potential
cost at most $3(r'(x)-r(x))$; a final single rotation adds at most one.

At a preferred-path switch, the total weight of a detached or attached
auxiliary path equals the represented subtree size of the demoted or
promoted child.  The change in the splayed root's own weight compensates
exactly, so the switch does not change $\Phi$.  Along an upward access,
the next node's initial auxiliary weight contains the preceding exposed
path through its dashed-child weight.  The rank differences therefore
telescope to $O(\log N_f)$, apart from a constant per access iteration.

To pay for those iterations, call a represented edge heavy if its
child's subtree has more than half its parent's size, and add a
sufficiently scaled count of dashed heavy edges to the potential.
A represented root path has at most $\log_2 N_f$ light edges.  Access
promotes every dashed edge it traverses.  A heavy off-path edge is
demoted only opposite a light preferred child on the accessed path, or
at its endpoint.  The iteration count plus the change in this additional
potential is consequently $O(\log N_f)$.

It remains to check represented-tree changes.  After access, write the
weight of a path vertex as one plus the sizes of its off-path dashed
subtrees.  These weights are unchanged by reversing the complete
exposed path.  Recursive left-right swapping in its auxiliary splay tree
preserves each auxiliary subtree's vertex set and total weight, hence
preserves $\Phi$.  If an off-path child becomes heavy after reversal,
the next edge down the reversed path must be light, or the vertex is
the path endpoint.  There are only $O(\log N_f)$ such events.  Cutting
an accessed parent edge similarly leaves the path weights unchanged,
does not increase $\Phi$, and charges newly heavy off-path children to
light edges of the shortened path.  Linking a root below an accessed
vertex increases only that auxiliary root's weight sum; its rank grows
by at most $\log N_f$.  Existing off-path children of its ancestors
cannot newly become heavy as ancestor subtrees grow, and the new dashed
edge contributes at most one.  All potentials are nonnegative and the
initial singleton state has zero potential.  This proves the bound,
also under recycling of a fixed node pool.

For threshold reporting, expose the updated path, read a witnessed
nonpositive minimum, access that node, and disable it at a point.
Repeating these operations costs the claimed amount per actual disable.
\end{proof}

\subsection{Preserving exact values under forest changes}
\label{engine:exact-values}

Set
\begin{equation}\label{engine:copy-bound}
                              D=M^{10/a}.
\end{equation}
There are at most $D$ current forest copies of any actual edge.  Indeed,
the two endpoint occurrence sets have at most $D/2$ vertices together,
and the copies between them form a forest.  For a loop the same argument
uses just one occurrence set.  The constant slack in the vertex
multiplicity bound is more than sufficient.

Maintain a list of current copies for each actual edge.  For each real
coordinate keep an explicit signed base, with invariant
\begin{equation}\label{engine:value-invariant}
                   f_e=\mathrm{base}_e+
                           \sum_{c\text{ a current copy of }e} h_c.
\end{equation}
An explicit single-edge traversal changes the base.  Before deleting a
copy, read its $h_c$ and flush it into the base; a new copy starts with
$h_c=0$.  This procedure is used also on price replacements and
retargetings, with old mapped data retained until all removals have been
processed.  A client value setting to $f_e^{\rm new}$ sets the base to
$f_e^{\rm new}-\sum_c h_c$.  Consequently a point read, point setting,
or sweep of all copies costs $O(D)$ tree point operations, plus
dictionary work.

The invariant is independent of the copy's present role in the backend.
In particular, the public-forest construction in
Section~\ref{sec:backend-updates} may reassign a dormant arm slot to a
different child occurrence at the same label.  Its internal mapped edges
retain their lifetimes and physical images.  Their $h_c$ are therefore
retained, not restarted.  A previous increment need not be a circulation
in the current representation of the forest: its already performed
physical increment is recorded by~\eqref{engine:value-invariant}.
Every \emph{new} step is a circulation in the current mapped graph.
Thus representation maintenance alone neither changes a value nor loses
an earlier increment.

\subsection{Drift detection and refresh accounting}\label{engine:refresh}

Absolute counters have a separate purpose.  At an anchor of edge $e$,
reset all its current copy counters $A_c$ and one additional bucket.
This bucket stores absolute motion of explicit single-edge traversals
and absolute counts flushed from copies that subsequently die.  Choose
a power-of-two threshold $\theta_e$ with
\begin{equation}\label{engine:threshold}
          \frac{qH}{800D\ell_e}\le\theta_e
                          \le\frac{qH}{400D\ell_e}.
\end{equation}
The interval has ratio two, so it contains such a power.  Each copy
and the bucket use this same threshold.  A new copy starts
with $A_c=0$ and the current threshold.  Copy death transfers $A_c$ to
the bucket before removal, as well as transferring $h_c$ to the base.

On an updated path, report every enabled copy reaching threshold using
the witnessed minimum, and disable it pending service.  A disabled copy
continues to accumulate both signed and absolute motion.  The whole
circulation update is completed before any client change or reprice is
serviced.  An explicit traversal or a bucket flush is tested directly.
Each threshold event sets a coalesced pending flag on the \emph{variable},
not just on the copy.  The flag survives deletion of the copy that
caused it.

Service a pending variable by an exact point read and anchor/price
refresh, with a notification to the client.  The ordering is important:
first reset every current absolute counter and the bucket, reenable all
current copies, and clear the old pending flag; then process the backend
replacement and its forest logs.  The current register, prices, and
thresholds are installed consistently with this no-motion reset.  Its
own old-copy removals therefore flush zero absolute mass for that
coordinate.  Other coordinates flagged by those logs are queued for
later service, after the current transaction finishes.  Explicit client
settings and descriptor changes reanchor in the same way and either
service or subsume an old pending flag.  Stale queue entries are skipped
or coalesced.

All old mapped edge data remain available through removals.  New copies
receive their current mapped prices and thresholds before a path can be
used again.  Drain pending engine flags and return the register-update
list; the client can then settle its local changes before asking for
another pricing step.  Further notifications arising in that settlement
are handled by the same rule.  No motion occurs during these
transactions.

\begin{lemma}[Register accuracy and notification bound]
\label{engine:refresh-bound}
After settlement, \eqref{engine:register-guarantee} holds.  The total
number of engine refresh reports and copy-disable events is $O(D)$ times
the number of successful pricing steps, apart from explicit client
operations.  The assertion is unaffected by arbitrary forest churn,
dormant-slot reuse, or client reanchoring.
\end{lemma}

\begin{proof}
The signed displacement since the most recent anchor is bounded by the
sum of current absolute copy counters and the bucket.  This is true
initially, preserved by path updates and explicit traversals, and
preserved by transfer on copy death.  Dormant retained copies keep their
counters.  At a settled state there are at most $D$ copies and one
bucket, each below threshold.  Since
$\lambda_e(\widehat f_e)\le\ell_e/H$, we obtain
\[
 \lambda_e(\widehat f_e)|f_e-\widehat f_e|
 \le \frac{\ell_e}{H}(D+1)\theta_e
 \le \frac{q(D+1)}{400D}\le q/200.
\]
This is stronger than~\eqref{engine:register-guarantee}.

For the cost bound, partition raw real-edge motion into intervals between
successive anchors of each coordinate.  Its price is fixed on such an
interval.  An engine report consumes at least
\[
                       \theta_e\ell_e/H\ge q/(800D)
\]
of that interval's raw priced absolute motion.  The consumed intervals
are disjoint for one coordinate, and different coordinates consume
different real-edge motion.  Absolute mass moves between copies and the
bucket without duplication.  It is reset before price-replacement logs,
so it cannot be charged again at a new price.  A client anchor may
discard old drift credit, but does not create new credit and is charged
as an external operation.

Similarly, a disable event can charge the motion of its copy since
creation or its latest reset.  A copy is disabled at most once during
that uninterrupted accumulation; reenabling at an anchor resets it.
These disable charges are disjoint within this second charge category.
Counting both report charges and disable charges only introduces a
constant factor.  By~\eqref{engine:step-scale}, a successful step creates
at most $q/200$ of raw priced absolute motion.  The claimed $O(D)$
amortized bound follows.

Finally, a no-motion refresh cascade is finite.  A queued variable can
consume its extant absolute mass once; after its reset, all its absolute
locations are zero and remain so until new motion.  Its own replacement
logs cannot recreate that mass.  Other variables may expose previously
accumulated mass, but their reports have the disjoint charges just
described.  Pending flags are per variable, so copy deletion loses
neither the required service nor its charge.  This proves both
termination of settlement and the notification bound.
\end{proof}

\subsection{Precision, total cost, and completion of the engine proof}
\label{engine:implementation}

All quantities maintained by the engine can be represented exactly with
charged multiword arithmetic.  The only nonintegral input values and
increments are on the common dyadic grid.  In particular, a positive
logarithmic argument is $k\varepsilon$ for an integer $k>0$ with
$O(\log S)$ bits, since the argument is at most $S^{20}$ and
$\varepsilon=S^{-3000}$.  Its reciprocal is
$\varepsilon^{-1}/k$.  Each price uses at most two such fractions and
one bounded integer weight, so their products, sums, and quotient
operations still have $O(\log S)$-bit inputs and outputs.
Rational comparisons of scores and scale choices likewise use bounded
integer operations.  No barrier value or ideal real-valued probability
is evaluated.

For a concrete precision check, bound the total number of traversals in
a returned representation by $S^{20}$, using its piece count and the
forest-size bound rather than expanding its paths.  The check holds
eventually for each fixed sufficiently large parameter, since
$M\le S^2$ and both piece count and occurrence overhead have exponent
$O(1/a)+o_a(1)$.  Equations~\eqref{engine:derivative-range} and
\eqref{engine:integer-prices} then give the loose bound
\begin{equation}\label{engine:precision-bounds}
                 V/H\le S^{1044},\qquad
                 |\mu|\ge S^{-1045}/400.
\end{equation}
For sufficiently large checked $S$, a power of two satisfying
\eqref{engine:step-scale} is therefore an exact multiple of
$\varepsilon=S^{-3000}$.  Each piece has integral signed traversal
coefficients, so every partial update and the final circulation are on
the grid, with no rounding of conservation equations.  Thresholds fit
the same grid: even using $D\le S^2$ gives
$\theta_e\ge S^{-1025}/800$.  These inequalities can be checked
directly, with suspension at inadequate sizes.

Since $V\ge1$, step scalars have magnitude at most $S^{100}$.
There are at most $Q\le S^{100}$ successful steps, each with at most
$S^{20}$ raw traversals.  Thus their total unweighted absolute motion is
at most $S^{220}$.  Transfers between copies and bases or buckets do not
duplicate this motion; anchors reset absolute counters, and client
settings replace values of magnitude at most $S^{20}$.  Signed and
absolute tags, sums, buckets, and bases consequently have polynomial
magnitudes on the common grid.  Their bit lengths are $O_a(\log S)$ at
fixed parameters; exact pricing and key operations add only
logarithmic-power overhead.  The backend separately accounts
for its arithmetic and imported static flow calls.  As in the final
uniform implementation, larger logical integers in branches with varying
parameters are digit arrays whose construction and operations are
incrementally charged, not atomic unbounded-word instructions.

Reserve and recycle a pool of
$N_f=O(M^{1+10/a})$ subdivided forest nodes.  The initial pool and forest
cost fit $M^{1+O(1/a)+o_a(1)}$.  By
Lemma~\ref{engine:tree-cost}, processing each exposed forest-log record
or path piece costs logarithmic amortized overhead.  A coordinate read,
reset, or price change uses at most $O(D)$ point operations for its copy
list.  By Lemma~\ref{engine:refresh-bound}, successful steps cause only
$O(D)$ reports and disables amortized.  Repricing those reported
coordinates invokes only that many backend replacements, each with the
cost and explicit recourse of Lemma~\ref{lem:flat-cycle}.  Thus the total
of backend work, forest-log processing, path operations, copy sweeps,
and threshold service is
\[
                   (M+Q)M^{O(1/a)+o_a(1)}.
\]
Using the chosen rounding of $M$ and
\eqref{engine:nominal-size} gives~\eqref{engine:time}.  A
\textsc{Low} attempt uses the same query overhead and no scan of all
coordinates.  Explicit changes and all notifications are settled before
the next query; no hidden full-vector update is involved.

\begin{proof}[Proof of Lemma~\ref{lem:priced-steps}]
Apply Lemma~\ref{lem:flat-cycle} to the port graph with prices
\eqref{engine:integer-prices}.  Exact point operations and forest updates
preserve~\eqref{engine:value-invariant}.  The refresh procedure proves
\eqref{engine:register-guarantee}, and therefore validates
\eqref{engine:price-accuracy} before each pricing call.  The lifting,
quality margin, and acceptance test prove part~\textup{(i)}.
Equations~\eqref{engine:accepted-direction}--\eqref{engine:step-scale},
the partial-motion bound, and the grid check prove part~\textup{(ii)},
including exact circulation and positivity.  The objective estimate
following the statement proves its progress assertion.
Lemma~\ref{engine:refresh-bound} proves part~\textup{(iii)}, and the
preceding cost accounting proves part~\textup{(iv)}.

Backend failures are explicitly propagated.  Its conditional instance
budget gives $2^{-\sigma}S^{-60}$ for the whole engine instance; the
remaining operations are deterministic.  Representation, anchor-range,
and precision prerequisites are checked on their local inputs, with
counts maintained incrementally.  During a whole step, the already
proved raw-motion bound supplies the required relative margins; no
global scan is needed to certify them.  The operation and record guards
are applied to prefixes, and the preceding bounds hold through a
putative first guard-reaching operation.  For each sufficiently large
fixed $a$ and each fixed $\sigma\ge0$, these bounds fall below the loose polynomial caps, and
the size and quality inequalities hold for all sufficiently large $S$.
Together with the backend's guarded-prefix guarantees in
Section~\ref{sec:backend-updates}, this proves eventual adequacy on
every path, independently of whether a stochastic routine ultimately
reports failure.  Such a failure is propagated promptly and is never
reclassified as a deterministic suspension.
\end{proof}
